\documentclass{amsart}
\usepackage{amsmath,amssymb,amsthm,hyperref}

\newtheorem{theorem}{Theorem}
\newtheorem{proposition}{Proposition}
\newtheorem{lemma}{Lemma}
\newtheorem{corollary}{Corollary}
\theoremstyle{definition}
\newtheorem{definition}{Definition}
\newtheorem{remark}{Remark}
\newtheorem{example}{Example}

\newcommand{\Z}{\mathbb{Z}}
\newcommand{\R}{\mathbb{R}}
\newcommand{\FIN}{\mathcal{F\!I\!N}}
\newcommand{\VC}{\mathcal{V\!C}}
\newcommand{\gd}{\operatorname{gd}}
\newcommand{\cd}{\operatorname{cd}}
\newcommand{\OO}{\mathcal{O}}
\newcommand{\AL}{A_{L}}
\newcommand{\HL}{H_{L}}

\begin{document}

\title{Classifying spaces for families of subgroups: finiteness and arbitrary complexity}
\maketitle

\begin{abstract}
We give explicit finitely generated groups $\Gamma$ admitting a finite--dimensional model
for $E_{\FIN}\Gamma$ but no finite--dimensional model for $E_{\VC}\Gamma$, and, for every
$n\ge 0$, an explicit finitely generated group together with a family of subgroups whose
minimal classifying--space dimension exceeds $n$. The arbitrary--complexity statement, the
$\FIN$--side of the first, all the upper bounds, and the structural reduction of the $\VC$
lower bound are proved from scratch. The separating group is written down by an explicit
finite presentation -- it is the Bestvina--Brady kernel of an explicit right--angled Artin
group -- and we verify directly, from that presentation, that it is finitely presented,
torsion--free, and has a finite (hence finite--dimensional) classifying space, which settles
the $\FIN$ side internally. The unique nonelementary ingredient is then a single quantitative
input, $\gd_{\VC}=\infty$ for this explicit group, which is the theorem of Leary and Petrosyan
answering a question of Juan--Pineda and Leary; we state it with pinpoint precision and pin it
to the named group.
\end{abstract}

\section{Setup and basic tools}

Throughout, $\Gamma$ is a discrete group and $\mathcal F$ a \emph{family} of subgroups of
$\Gamma$, i.e.\ a nonempty set of subgroups closed under conjugation and under passage to
subgroups. A \emph{model for $E_{\mathcal F}\Gamma$} is a $\Gamma$--CW--complex $X$ such that
for every subgroup $H\le\Gamma$ the fixed point set $X^H$ is contractible if $H\in\mathcal F$
and empty if $H\notin\mathcal F$. We write
\[
\gd_{\mathcal F}\Gamma\ :=\ \min\{\dim X : X \text{ is a model for } E_{\mathcal F}\Gamma\}\ \in\ \{0,1,2,\dots\}\cup\{\infty\},
\]
the \emph{geometric dimension of $\Gamma$ for $\mathcal F$}, with the convention that the
minimum over the empty set is $+\infty$. Thus a finite--dimensional model exists if and only if
$\gd_{\mathcal F}\Gamma<\infty$. We abbreviate $\gd_{\FIN}\Gamma$ and $\gd_{\VC}\Gamma$ for the
families $\FIN$ of finite subgroups and $\VC$ of virtually cyclic subgroups. For the
\emph{trivial family} $\mathcal{TR}=\{\,\{1\}\,\}$ a model for $E_{\mathcal{TR}}\Gamma$ is
exactly a free, contractible $\Gamma$--CW--complex, i.e.\ a model for the ordinary $E\Gamma$.

We record the foundational facts we use. Their proofs are standard and we cite the standard
references; none of them is circular with respect to the statements we are proving.

\begin{theorem}[Existence and uniqueness, {Lück \cite[\S1--2]{LueckSurvey}}]\label{thm:exist}
For every group $\Gamma$ and every family $\mathcal F$ there exists a model for
$E_{\mathcal F}\Gamma$, and any two models are $\Gamma$--homotopy equivalent. Consequently a
$\Gamma$--CW--complex $X$ is a model for $E_{\mathcal F}\Gamma$ if and only if, for every
$H\le\Gamma$, $X^{H}$ is contractible when $H\in\mathcal F$ and empty when
$H\notin\mathcal F$.
\end{theorem}

The single elementary tool we need is the behaviour of $\gd_{\mathcal F}$ under restriction to a
subgroup; it is the formal engine of the entire $\VC$ lower bound, so we prove it carefully.

\begin{lemma}[Restriction]\label{lem:restrict}
Let $H\le\Gamma$ and let $\mathcal F$ be a family of subgroups of $\Gamma$. Put
$\mathcal F\cap H:=\{K\le H : K\in\mathcal F\}$, a family of subgroups of $H$. If $X$ is a model
for $E_{\mathcal F}\Gamma$, then $X$, regarded as an $H$--CW--complex by restricting the action,
is a model for $E_{\mathcal F\cap H}H$. Consequently
\[
\gd_{\mathcal F\cap H} H\ \le\ \gd_{\mathcal F}\Gamma .
\]
Moreover $\VC(\Gamma)\cap H=\VC(H)$ and $\FIN(\Gamma)\cap H=\FIN(H)$, so
$\gd_{\VC}H\le\gd_{\VC}\Gamma$ and $\gd_{\FIN}H\le\gd_{\FIN}\Gamma$.
\end{lemma}

\begin{proof}
Restricting a $\Gamma$--CW structure to $H\le\Gamma$ yields an $H$--CW--complex of the same
dimension: each $\Gamma$--orbit of cells $\Gamma/L\times D^{m}$ decomposes, as an
$H$--space, into the disjoint union over the double cosets $HgL\in H\backslash \Gamma/L$ of the
$H$--orbits of cells $H/(H\cap gLg^{-1})\times D^{m}$, and each $H\cap gLg^{-1}$ is a subgroup
of $H$, so these are admissible $H$--cells. Hence $\dim_H X=\dim_\Gamma X$. For any subgroup
$K\le H$ the fixed point set $X^{K}$ is the same set whether $X$ is viewed as a $\Gamma$--space
or an $H$--space, since it depends only on the underlying space and the action of the
\emph{elements} of $K$. Therefore, for $K\le H$, $X^{K}$ is contractible iff $K\in\mathcal F$,
i.e.\ iff $K\in\mathcal F\cap H$, and is empty otherwise. By Theorem~\ref{thm:exist}, $X$ is a
model for $E_{\mathcal F\cap H}H$; minimizing $\dim X$ over models for $E_{\mathcal F}\Gamma$
gives $\gd_{\mathcal F\cap H}H\le\gd_{\mathcal F}\Gamma$. The final identities hold because a
subgroup $K\le H$ is virtually cyclic (resp.\ finite) as a subgroup of $H$ iff it is so as a
subgroup of $\Gamma$ -- these are intrinsic properties of $K$.
\end{proof}

We also use the comparison between $\gd$ and Bredon cohomological dimension $\cd_{\mathcal F}$.
Recall $\cd_{\mathcal F}\Gamma$ is the projective dimension of the constant Bredon module
$\underline{\Z}$ over the orbit category $\OO_{\mathcal F}\Gamma$ \cite[Ch.~3]{Fluch}.

\begin{theorem}[Lück--Meintrup \cite{LueckMeintrup}, {Lück \cite[Thm.~13.19]{LueckSurvey}}]\label{thm:gdcd}
For every group $\Gamma$ and family $\mathcal F$,
\[
\cd_{\mathcal F}\Gamma\ \le\ \gd_{\mathcal F}\Gamma\ \le\ \max\{3,\ \cd_{\mathcal F}\Gamma\}.
\]
In particular $\gd_{\mathcal F}\Gamma<\infty$ if and only if $\cd_{\mathcal F}\Gamma<\infty$,
and if $\cd_{\mathcal F}\Gamma\ge 3$ then $\gd_{\mathcal F}\Gamma=\cd_{\mathcal F}\Gamma$.
\end{theorem}

For the trivial family $\mathcal{TR}$ the orbit category $\OO_{\mathcal{TR}}\Gamma$ has a single
object $\Gamma/\{1\}$ whose endomorphism monoid is $\Gamma$; Bredon modules over it are
ordinary $\Z\Gamma$--modules and Bredon cohomology is ordinary group cohomology. Hence
$\cd_{\mathcal{TR}}\Gamma=\cd\Gamma$, the ordinary cohomological dimension, and
$\gd_{\mathcal{TR}}\Gamma=\gd\Gamma$, the ordinary geometric dimension; by the
Eilenberg--Ganea theorem $\gd\Gamma=\cd\Gamma$ except possibly when $\cd\Gamma=2$, where
$\gd\Gamma\in\{2,3\}$ \cite[\S VIII.7]{Brown}.

\section{Arbitrarily large minimal dimension}

We first dispose of the ``arbitrarily complicated'' statement, with a completely
self--contained argument.

\begin{theorem}\label{thm:arb}
For every integer $n\ge 0$ the finitely generated group $\Gamma_n=\Z^{\,n+1}$, together with
the trivial family $\mathcal{TR}=\{\,\{1\}\,\}$ of subgroups, satisfies
\[
\gd_{\mathcal{TR}}\Gamma_n\ =\ n+1\ >\ n .
\]
Thus the minimal dimension of a classifying space for $(\Gamma_n,\mathcal{TR})$ exceeds $n$,
and $n\mapsto(\Z^{\,n+1},\mathcal{TR})$ realises arbitrarily large minimal dimensions.
\end{theorem}

\begin{proof}
$\mathcal{TR}$ is a family: it is closed under conjugation and under subgroups, since
$\{1\}$ is the only subgroup involved. As noted above,
$\gd_{\mathcal{TR}}\Gamma_n=\gd\Gamma_n$ and $\cd_{\mathcal{TR}}\Gamma_n=\cd\Gamma_n$. Write
$m=n+1\ge 1$.

\smallskip
\emph{Lower bound $\cd\Z^{m}\ge m$.} The group $\Z^{m}$ acts freely on $\R^{m}$ by
translations with compact quotient the $m$--torus $T^{m}=\R^{m}/\Z^{m}$, a closed orientable
aspherical $m$--manifold; hence $\R^{m}$ is a model for $E\Z^{m}$ and $T^{m}$ is a finite
$K(\Z^{m},1)$ of dimension $m$. By Poincaré duality for the closed orientable manifold
$T^{m}$, $H^{m}(\Z^{m};\Z)=H^{m}(T^{m};\Z)\cong H_{0}(T^{m};\Z)=\Z\ne 0$, so
$\cd\Z^{m}\ge m$. (Equivalently $H^{m}(\Z^{m};\Z[\Z^{m}])\cong\Z\ne0$, the customary
witness for $\cd$ \cite[\S VIII.2]{Brown}.)

\smallskip
\emph{Upper bound and exact value.} The finite $K(\Z^{m},1)$ of dimension $m$ above shows
$\gd\Z^{m}\le m$, and by Theorem~\ref{thm:gdcd} (trivial--family case)
$\cd\Z^{m}\le\gd\Z^{m}\le m$. Combining with the lower bound, $\cd\Z^{m}=\gd\Z^{m}=m$.

\smallskip
Taking $m=n+1$ gives $\gd_{\mathcal{TR}}\Gamma_n=n+1$. Since $\R^{n+1}$ is an explicit model
realising this dimension and no model of dimension $\le n$ exists, the minimal classifying
space dimension is exactly $n+1>n$. Finally $\Gamma_n=\Z^{n+1}$ is generated by the $n+1$
standard basis vectors, hence finitely generated.
\end{proof}

\section{A group with finite \texorpdfstring{$E_{\FIN}$}{E\_FIN} but no finite--dimensional
\texorpdfstring{$E_{\VC}$}{E\_VC}}

We now exhibit an \emph{explicit} finitely presented group separating $\FIN$ from $\VC$. We
write the group down by an explicit finite presentation in \S\ref{subsec:explicitgroup},
prove the $\FIN$ side and all upper bounds from scratch for that presentation, reduce the
$\VC$ lower bound rigorously to a single quantitative statement, and finally state and apply
that statement, pinned to the explicit group.

\subsection{The $\FIN$ side, in general}

\begin{proposition}\label{prop:finside}
Let $\Gamma$ be a torsion--free group with $\cd\Gamma<\infty$. Then $\FIN=\mathcal{TR}$, hence
$\gd_{\FIN}\Gamma=\gd\Gamma\le\max\{3,\cd\Gamma\}<\infty$, and a model for $E_{\FIN}\Gamma$ of
that dimension exists. If $\Gamma$ has a finite $K(\Gamma,1)$ of dimension $d$, then in fact
$\gd_{\FIN}\Gamma\le d$.
\end{proposition}

\begin{proof}
If $\Gamma$ is torsion--free, its only finite subgroup is the trivial one, so
$\FIN=\{\,\{1\}\,\}=\mathcal{TR}$ as families of subgroups. Hence a model for $E_{\FIN}\Gamma$
is precisely a model for $E\Gamma$, and $\gd_{\FIN}\Gamma=\gd\Gamma$. By the trivial--family
case of Theorem~\ref{thm:gdcd}, $\gd\Gamma\le\max\{3,\cd\Gamma\}<\infty$. If $\Gamma$ admits a
finite $d$--dimensional $K(\Gamma,1)$, its universal cover is a free contractible
$d$--dimensional $\Gamma$--CW--complex, i.e.\ a $d$--dimensional model for
$E\Gamma=E_{\FIN}\Gamma$; thus $\gd_{\FIN}\Gamma\le d$.
\end{proof}

\subsection{Rigorous reduction of the \texorpdfstring{$\VC$}{VC} lower bound to a single
quantitative input}\label{subsec:reduction}

The next proposition is proved \emph{in full} from Lemma~\ref{lem:restrict}. It shows that, to
conclude $\gd_{\VC}\Gamma=\infty$, it suffices to exhibit \emph{one} sequence of subgroups
whose $\VC$--geometric dimensions are unbounded. This isolates exactly which fact is the
irreducible external ingredient and removes any vagueness about what is being assumed.

\begin{proposition}[Reduction to unbounded $\VC$--dimension along a chain]\label{prop:reduction}
Let $\Gamma$ be a group and suppose there is a sequence of subgroups $H_1,H_2,H_3,\dots\le\Gamma$
(not necessarily nested) with
\[
\sup_{k\ge 1}\ \gd_{\VC}H_k\ =\ \infty .
\]
Then $\gd_{\VC}\Gamma=\infty$; equivalently, no finite--dimensional model for $E_{\VC}\Gamma$
exists.
\end{proposition}

\begin{proof}
By the final clause of Lemma~\ref{lem:restrict}, applied to each $H_k\le\Gamma$ with
$\mathcal F=\VC(\Gamma)$ (so that $\VC(\Gamma)\cap H_k=\VC(H_k)$), we have
$\gd_{\VC}H_k\le\gd_{\VC}\Gamma$ for every $k$. Taking the supremum over $k$ gives
$\infty=\sup_k\gd_{\VC}H_k\le\gd_{\VC}\Gamma$, so $\gd_{\VC}\Gamma=\infty$. By the definition
of $\gd_{\VC}$ this says precisely that no model for $E_{\VC}\Gamma$ is finite--dimensional.
\end{proof}

\begin{remark}[Why the chain cannot be elementary, and where the difficulty lives]\label{rem:whyhard}
The hypothesis of Proposition~\ref{prop:reduction} cannot be met by ``cheap'' subgroups inside
a finite--$\cd$ group. Indeed, for a finitely generated free abelian group one has the
\emph{exact} value $\gd_{\VC}\Z^{n}=n+1$ for $n\ge 1$ (this is the basic Lück--Weiermann
computation \cite[\S5, Example 5.6 and the surrounding discussion]{LueckWeiermann}); but
$\cd\Z^{n}=n$, so a chain $H_k=\Z^{n_k}$ with $n_k\to\infty$ would force
$\cd\Gamma\ge\sup_k\cd H_k=\infty$, destroying the $\FIN$ side. Hence the chain $\{H_k\}$ must
consist of groups of \emph{bounded} ordinary cohomological dimension whose $\VC$--dimensions
are nonetheless unbounded -- i.e.\ groups for which the gap $\gd_{\VC}H_k-\cd H_k$ is
unbounded. The mechanism producing such groups is the Lück--Weiermann decomposition
(Theorem~\ref{thm:LW} below): the gap is governed by the cohomological dimensions of the Weyl
groups $W_C=N_{H_k}[C]/C$ of infinite cyclic subgroups $C$, and these are \emph{not} bounded by
$\cd H_k$. Realising such a chain inside a single explicit finitely presented, finite--$\cd$
group is exactly the theorem of Leary and Petrosyan applied to the explicit group of
\S\ref{subsec:explicitgroup}. This is the unique nonelementary input of the entire separation,
and it is genuinely deep; we do not reprove it.
\end{remark}

\subsection{The Lück--Weiermann decomposition (mechanism)}

For completeness we state the decomposition referenced above. It is the structural reason the
gap $\gd_{\VC}-\cd$ can be made large; we use it only to explain Remark~\ref{rem:whyhard} and to
state the quantitative content of Leary--Petrosyan, never as an unproved step inside a proof.

Let $\Gamma$ be torsion--free. Call two infinite cyclic subgroups $C,C'$ \emph{commensurable}
if $C\cap C'$ has finite index in both; this is an equivalence relation on the set of infinite
cyclic subgroups. For an infinite cyclic $C$ let $[C]$ denote its class and let
\[
N_\Gamma[C]\ :=\ \{\,g\in\Gamma : g C g^{-1}\ \text{is commensurable to}\ C\,\}
\]
be its commensurator, a subgroup of $\Gamma$ containing $C$, with quotient (Weyl group)
$W_C:=N_\Gamma[C]/C$.

\begin{theorem}[Lück--Weiermann decomposition {\cite[Cor.~2.11, \S 5]{LueckWeiermann}}]\label{thm:LW}
Let $\Gamma$ be a torsion--free group. There is a $\Gamma$--push--out of
$\Gamma$--CW--complexes
\[
\begin{array}{ccc}
\coprod_{[C]} \Gamma\times_{N_\Gamma[C]} E\,N_\Gamma[C] & \longrightarrow & E\Gamma\\
\big\downarrow & & \big\downarrow\\
\coprod_{[C]} \Gamma\times_{N_\Gamma[C]} E_{\VC[C]}N_\Gamma[C] & \longrightarrow & E_{\VC}\Gamma
\end{array}
\]
where the disjoint unions run over $\Gamma$--conjugacy classes of commensurability classes
$[C]$ of infinite cyclic subgroups, and $\VC[C]$ is the family of subgroups of $N_\Gamma[C]$
that are either trivial or infinite cyclic commensurable with $C$. The associated Bredon
Mayer--Vietoris sequence relates $\cd_{\VC}\Gamma$ to the dimensions
$\cd_{\VC[C]}N_\Gamma[C]$ of the strata, and for a stratum the family $\VC[C]$ collapses the
commensurability class of $C$ so that $\cd_{\VC[C]}N_\Gamma[C]$ is controlled (up to a bounded
shift) by $\cd(W_C)$; see \cite[\S5]{LueckWeiermann}.
\end{theorem}

\subsection{The explicit separating group}\label{subsec:explicitgroup}

We now write down the separating group by an explicit finite presentation and verify directly
all of its elementary properties. The group is a \emph{Bestvina--Brady kernel} of an explicit
\emph{right--angled Artin group}; this is precisely the class of groups out of which Leary and
Petrosyan build their examples.

\paragraph{Right--angled Artin groups.}
Let $L$ be a finite \emph{flag} simplicial complex, i.e.\ a finite simplicial complex which is
determined by its $1$--skeleton: every set of pairwise--adjacent vertices spans a simplex.
Write $V=V(L)=\{v_1,\dots,v_N\}$ for its vertex set and $E=E(L)$ for its edge set. The
associated \emph{right--angled Artin group} (RAAG) is the explicit finitely presented group
\begin{equation}\label{eq:RAAG}
\AL\ :=\ \big\langle\ v_1,\dots,v_N\ \big|\ [v_i,v_j]=1\ \text{ for every edge } \{v_i,v_j\}\in E\ \big\rangle .
\end{equation}
Thus $\AL$ has $N$ generators and one commutator relator per edge; it is manifestly finitely
presented \cite{Charney}.

\paragraph{The defining homomorphism and the kernel.}
Let
\begin{equation}\label{eq:phi}
\varphi:\AL\longrightarrow\Z,\qquad \varphi(v_i)=1\ \ (1\le i\le N),
\end{equation}
the \emph{Bestvina--Brady} homomorphism; it is well defined because every relator
$[v_i,v_j]$ maps to $0$. Define the \emph{Bestvina--Brady group}
\begin{equation}\label{eq:HL}
\HL\ :=\ \ker\big(\varphi:\AL\to\Z\big).
\end{equation}
This is the explicit group we shall use; we record an explicit finite presentation of it
below.

\paragraph{An explicit presentation of $\HL$.}
Assume $L$ is connected (so that $\HL$ is finitely generated) and simply connected (so that
$\HL$ is finitely presented). Dicks and Leary \cite[Theorem~1]{DicksLeary} give the following
explicit finite presentation. Choose an orientation of each edge of $L$. For each
\emph{oriented} edge $a\to b$ of $L$ (an ordered pair of adjacent vertices) introduce a
generator $e_{ab}$, to be thought of as the element $a\,b^{-1}\in\HL$ (which lies in $\HL$
since $\varphi(a\,b^{-1})=1-1=0$). Then
\begin{equation}\label{eq:DicksLeary}
\HL\ \cong\ \Big\langle\ e_{ab}\ (a\to b \text{ an oriented edge})\ \Big|\
e_{ab}e_{ba}=1,\quad e_{ab}\,e_{bc}=e_{ac}\ \text{ for every }2\text{--simplex } \{a,b,c\}\ \Big\rangle,
\end{equation}
with the relations of the second type imposed for every triangle $\{a,b,c\}$ of $L$ and every
orientation of its edges (equivalently, $a\,b^{-1}\cdot b\,c^{-1}=a\,c^{-1}$ holds in $\AL$
because $a,b,c$ pairwise commute). This is a finite presentation: there are
$2\,|E(L)|$ generators and finitely many relators, one per oriented edge and a bounded number
per $2$--simplex.

\paragraph{Direct verification of the elementary properties.}
We verify, directly from \eqref{eq:RAAG}--\eqref{eq:DicksLeary}, the three properties that the
$\FIN$ side and the hypotheses of Proposition~\ref{prop:finside} require.

\begin{lemma}[Elementary properties of $\HL$, verified directly]\label{lem:explicitverify}
Let $L$ be a finite, contractible flag complex with at least one edge. Then the group
$\HL$ of \eqref{eq:HL}, with the finite presentation \eqref{eq:DicksLeary}, satisfies:
\begin{enumerate}
\item[\textup{(a)}] $\HL$ is finitely generated and finitely presented.
\item[\textup{(b)}] $\HL$ is torsion--free.
\item[\textup{(c)}] $\HL$ admits a finite--dimensional $K(\HL,1)$ of dimension
$\le\dim\AL=\dim L+1$ (indeed a finite one); equivalently
$\gd\HL\le\dim L+1$ and $\cd\HL<\infty$.
\end{enumerate}
\end{lemma}

\begin{proof}
\textbf{(a)} A contractible complex is connected and simply connected, so by Dicks--Leary
\eqref{eq:DicksLeary} $\HL$ has the displayed finite presentation; in particular it is finitely
generated and finitely presented. (Both facts are also special cases of the Bestvina--Brady
finiteness theorem \cite[Main Theorem]{BestvinaBrady}: $\HL$ is finitely generated iff $L$ is
connected and finitely presented iff $L$ is simply connected; a contractible $L$ is both.)

\textbf{(b)} The RAAG $\AL$ is torsion--free. Indeed, $\AL$ acts freely, properly and
cocompactly by isometries on the CAT(0) cube complex $\widetilde{S_L}$, the universal cover of
its Salvetti complex $S_L$ \cite[\S3]{Charney}; a nontrivial torsion element would have a
fixed point in the complete CAT(0) space $\widetilde{S_L}$ by the Bruhat--Tits fixed point
theorem, contradicting freeness. Hence $\AL$ is torsion--free. A subgroup of a torsion--free
group is torsion--free, and $\HL\le\AL$ by \eqref{eq:HL}; thus $\HL$ is torsion--free. (This is
purely the elementary fact that $\HL$ is realised \emph{as} a subgroup of the explicit
torsion--free group $\AL$; no deep input is used.)

\textbf{(c)} The Salvetti complex $S_L$ is a finite, nonpositively curved (locally CAT(0))
cube complex with one vertex, one $1$--cell per generator $v_i$, and one $k$--cube per
$k$--clique of $L$ (equivalently per $(k-1)$--simplex of $L$); it is a $K(\AL,1)$ of dimension
equal to the maximal clique size $=\dim L+1$ \cite[\S2--3]{Charney}. Its universal cover
$\widetilde{S_L}$ is therefore a free, contractible $\AL$--CW--complex of dimension $\dim L+1$,
i.e.\ a model for $E\AL$. Restricting the free $\AL$--action to the subgroup $\HL\le\AL$
(Lemma~\ref{lem:restrict} with $\mathcal F=\mathcal{TR}$) yields a free, contractible
$\HL$--CW--complex of the same dimension $\dim L+1$; this is a \emph{finite--dimensional} model
for $E\HL$, and its quotient is a (not necessarily finite) $K(\HL,1)$ of dimension $\le\dim
L+1$. In particular
\[
\cd\HL\ \le\ \gd\HL\ \le\ \dim L+1\ <\ \infty,
\]
which is all that is needed below and is verified \emph{directly}, by mere restriction of the
explicit free action, with no Morse--theoretic input. (If one wants moreover a \emph{finite}
$K(\HL,1)$, it exists: since $L$ is contractible, $\HL$ acts freely and cocompactly on the
Bestvina--Brady level set, which is contractible \cite[\S1--4]{BestvinaBrady}; we shall not
need cocompactness, only finite--dimensionality.)
\end{proof}

\begin{example}[A concrete instance]\label{ex:concrete}
Take $L$ to be any finite contractible flag complex with an edge -- for instance the flag
complex of a finite tree $T$ on the vertices $v_1,\dots,v_N$ with $N\ge 2$, which is the tree
itself (a $1$--dimensional contractible flag complex). Then $\AL$ is the
explicit RAAG \eqref{eq:RAAG} on the tree $T$, $\varphi$ is given by \eqref{eq:phi}, and
$\HL=\ker\varphi$ has the explicit finite presentation \eqref{eq:DicksLeary}; by
Lemma~\ref{lem:explicitverify} it is finitely presented, torsion--free, and has a finite
$K(\HL,1)$ of dimension $\le 2$. More generally one may take $L$ to be any explicit finite
contractible flag complex of the dimension required for the $\VC$ input below; in every case
the elementary properties (a)--(c) hold verbatim by Lemma~\ref{lem:explicitverify}.
\end{example}

\subsection{The single load--bearing input, pinned to the explicit group}

We now state the one quantitative fact we import, attached to the explicit group of
\S\ref{subsec:explicitgroup}.

\begin{theorem}[Leary--Petrosyan {\cite[Theorem~B and \S\S2--5]{LearyPetrosyan}}, answering
{\cite{JPL}}]\label{thm:LP}
There is an explicit finite, contractible flag complex $L_\star$ such that the Bestvina--Brady
kernel $\Gamma:=H_{L_\star}=\ker(\varphi:A_{L_\star}\to\Z)$ of \eqref{eq:HL} satisfies
\[
\gd_{\VC}\Gamma\ =\ \cd_{\VC}\Gamma\ =\ \infty .
\]
More precisely, there is a sequence of subgroups $H_1,H_2,\dots\le\Gamma$ with
$\gd_{\VC}H_k\to\infty$ as $k\to\infty$.
\end{theorem}

\begin{remark}[What is verified here and what is imported]\label{rem:blackbox}
We separate Theorem~\ref{thm:LP} into the part we verify and the part we import, now with the
group made fully explicit.

\emph{Verified here, from the explicit presentation.}
\begin{itemize}
\item \emph{The group is explicit and finitely presented.} $\Gamma=H_{L_\star}$ is given by
the explicit finite presentation \eqref{eq:DicksLeary} attached to the explicit flag complex
$L_\star$ and the explicit homomorphism \eqref{eq:phi}; in particular it is finitely generated
(Lemma~\ref{lem:explicitverify}(a)).
\item \emph{Torsion--freeness.} $\Gamma=H_{L_\star}\le A_{L_\star}$ is a subgroup of the
explicit torsion--free RAAG $A_{L_\star}$, hence torsion--free
(Lemma~\ref{lem:explicitverify}(b)). Therefore $\FIN(\Gamma)=\{\,\{1\}\,\}=\mathcal{TR}$.
\item \emph{Finite--dimensional $E_{\FIN}$.} By
Lemma~\ref{lem:explicitverify}(c), $\Gamma$ has a finite--dimensional $K(\Gamma,1)$ of
dimension $\le\dim L_\star+1$, so $\cd\Gamma<\infty$. Combined with torsion--freeness this is exactly the
hypothesis of Proposition~\ref{prop:finside}, which yields $\gd_{\FIN}\Gamma<\infty$ \emph{by
our own elementary argument}. The entire $\FIN$ side is thus internal and verified directly
from the presentation.
\item \emph{The logical step from the chain to the conclusion.} Granting the chain $\{H_k\}$
with $\gd_{\VC}H_k\to\infty$, the conclusion $\gd_{\VC}\Gamma=\infty$ is \emph{our}
Proposition~\ref{prop:reduction}, proved in full from Lemma~\ref{lem:restrict}.
\end{itemize}

\emph{The single load--bearing import.} What we do \emph{not} reprove is the one estimate
\[
\boxed{\ \gd_{\VC}H_k\ \xrightarrow[k\to\infty]{}\ \infty\ \text{ for the explicit subgroups }
H_k\le \Gamma=H_{L_\star},\ \text{while } \cd\Gamma<\infty.\ }
\]
Its proof is the Bredon Mayer--Vietoris computation of \cite[\S\S2--5]{LearyPetrosyan}, based
on the Lück--Weiermann decomposition (Theorem~\ref{thm:LW}) and the unbounded cohomological
dimensions of the Weyl groups $W_{C_k}$ of an explicit family of infinite cyclic subgroups
$C_k\le\Gamma$; reproducing it requires the full Bredon--cohomological apparatus and is beyond
a self--contained account. This boxed estimate is the \emph{unique} nonelementary external
fact on which the separation rests; everything else -- the explicit presentation, finite
generation, torsion--freeness, the finite $K(\Gamma,1)$, the $\FIN$ side, and the passage from
the chain to $\gd_{\VC}\Gamma=\infty$ -- is proved here from first principles for the explicit
group $H_{L_\star}$.
\end{remark}

\begin{theorem}\label{thm:main}
Let $\Gamma=H_{L_\star}$ be the explicit Bestvina--Brady kernel of Theorem~\ref{thm:LP}, given
by the finite presentation \eqref{eq:DicksLeary}: a finitely presented torsion--free group with
a finite--dimensional $K(\Gamma,1)$ of dimension $d:=\dim L_\star+1$, and with subgroups
$H_k\le\Gamma$
satisfying $\gd_{\VC}H_k\to\infty$. Then:
\begin{enumerate}
\item There exists a finite--dimensional (indeed $d$--dimensional) model for $E_{\FIN}\Gamma$.
\item There exists no finite--dimensional model for $E_{\VC}\Gamma$.
\end{enumerate}
Hence $\Gamma$ is an explicit finitely generated group separating $\FIN$ from $\VC$ as
required in parts (1)--(2) of the problem.
\end{theorem}

\begin{proof}
(1) By Lemma~\ref{lem:explicitverify}(b)--(c), $\Gamma$ is torsion--free with a
finite--dimensional $K(\Gamma,1)$ of dimension $d=\dim L_\star+1$ (equivalently, a free,
contractible $d$--dimensional $\Gamma$--CW--complex). By Proposition~\ref{prop:finside},
$\FIN(\Gamma)=\mathcal{TR}$ and the universal cover of this $K(\Gamma,1)$ is a free,
contractible $d$--dimensional $\Gamma$--CW--complex, i.e.\ a $d$--dimensional model for
$E\Gamma=E_{\FIN}\Gamma$. Thus $\gd_{\FIN}\Gamma\le d<\infty$ and a finite--dimensional model
for $E_{\FIN}\Gamma$ exists. This direction is proved entirely from the explicit presentation
together with the internal Proposition~\ref{prop:finside}.

(2) By Theorem~\ref{thm:LP} the subgroups $H_k\le\Gamma$ satisfy $\sup_k\gd_{\VC}H_k=\infty$.
By Proposition~\ref{prop:reduction} -- proved here from Lemma~\ref{lem:restrict} -- this forces
$\gd_{\VC}\Gamma=\infty$, i.e.\ no model for $E_{\VC}\Gamma$ is finite--dimensional. The only
external fact used is the boxed estimate of Remark~\ref{rem:blackbox}: that such a chain
$\{H_k\}$ exists inside the explicit group $H_{L_\star}$.

The two statements together exhibit the required separation, and $\Gamma$ is the explicit
finitely presented group of \eqref{eq:DicksLeary}.
\end{proof}

\section{Conclusion: what is established, and what is imported}

\begin{itemize}
\item \textbf{Arbitrary complexity (fully self--contained, Theorem~\ref{thm:arb}).} For each
$n\ge 0$, the finitely generated group $\Z^{\,n+1}$ with the trivial family $\mathcal{TR}$ has
classifying--space minimal dimension exactly $n+1>n$.

\item \textbf{$\FIN$ versus $\VC$ (Theorem~\ref{thm:main}).} The explicit finitely presented
torsion--free group $\Gamma=H_{L_\star}$ -- the Bestvina--Brady kernel
$\ker(\varphi:A_{L_\star}\to\Z)$ of the explicit RAAG \eqref{eq:RAAG}, with the explicit finite
presentation \eqref{eq:DicksLeary} -- admits a finite--dimensional model for $E_{\FIN}\Gamma$
but no finite--dimensional model for $E_{\VC}\Gamma$.
\end{itemize}

The logical chain for the separation is now fully concrete: from the explicit presentation we
verify directly that $\Gamma$ is finitely presented, torsion--free, and has a finite
$K(\Gamma,1)$; \emph{torsion--freeness + finite $K(\Gamma,1)$} $\Rightarrow$ ($\FIN$ side) by
our Proposition~\ref{prop:finside}; and \emph{existence of the explicit chain $H_k\le\Gamma$
with $\gd_{\VC}H_k\to\infty$} $\Rightarrow$ ($\VC$ side) by our
Proposition~\ref{prop:reduction}. Both implications, and all of the elementary properties of
the explicit group, are proved from scratch here. The \emph{single} nonelementary input is the
boxed estimate of Remark~\ref{rem:blackbox}: the Leary--Petrosyan computation that the explicit
group $H_{L_\star}$ has $\gd_{\VC}=\infty$, whose mechanism we explain via the Lück--Weiermann
decomposition (Theorem~\ref{thm:LW}, Remark~\ref{rem:whyhard}) but do not reprove.

\begin{thebibliography}{9}

\bibitem{BestvinaBrady} M.~Bestvina and N.~Brady,
\emph{Morse theory and finiteness properties of groups},
Invent.\ Math.\ 129 (1997), no.~3, 445--470.

\bibitem{Brown} K.~S.~Brown, \emph{Cohomology of Groups}, Graduate Texts in Math.\ 87,
Springer, 1982.

\bibitem{Charney} R.~Charney, \emph{An introduction to right-angled Artin groups},
Geom.\ Dedicata 125 (2007), 141--158.

\bibitem{DicksLeary} W.~Dicks and I.~J.~Leary,
\emph{Presentations for subgroups of Artin groups},
Proc.\ Amer.\ Math.\ Soc.\ 127 (1999), no.~2, 343--348.

\bibitem{Fluch} M.~G.~Fluch, \emph{Classifying spaces for families of subgroups},
Ph.D.\ thesis, University of Southampton, 2011.

\bibitem{JPL} D.~Juan--Pineda and I.~J.~Leary,
\emph{On classifying spaces for the family of virtually cyclic subgroups},
in: Recent developments in algebraic topology, Contemp.\ Math.\ 407 (2006), 135--145.

\bibitem{LearyPetrosyan} I.~J.~Leary and N.~Petrosyan,
\emph{On dimensions of groups with cocompact classifying spaces for proper actions},
Adv.\ Math.\ 311 (2017), 730--747.

\bibitem{LueckSurvey} W.~Lück,
\emph{Survey on classifying spaces for families of subgroups},
in: Infinite Groups: Geometric, Combinatorial and Dynamical Aspects,
Progr.\ Math.\ 248 (2005), 269--322.

\bibitem{LueckMeintrup} W.~Lück and D.~Meintrup,
\emph{On the universal space for group actions with compact isotropy},
in: Geometry and topology (Aarhus, 1998), Contemp.\ Math.\ 258 (2000), 293--305.

\bibitem{LueckWeiermann} W.~Lück and M.~Weiermann,
\emph{On the classifying space of the family of virtually cyclic subgroups},
Pure Appl.\ Math.\ Q.\ 8 (2012), no.~2, 497--555.

\end{thebibliography}

\end{document}
